% TOP_PROBLEM: {"id": 1355, "title": "Jónsson Algebra on ℵ_ω", "category_id": 10, "category": {"id": 10, "name": "set_theory", "display_name": "Set Theory", "description": "Foundations of mathematics, infinite sets, and cardinality.", "slug": "set-theory", "order_index": 10, "created_at": "2026-07-31T15:26:25.671Z"}, "status": "open"}
% FIRST_POSTED: 2026-10-01
% ATTEMPT_STATUS: unresolved
% Imported from: unsolvedmath_additional_500.tex; UnsolvedMath ID 1355
% !TeX program = lualatex
% Compile twice: lualatex 1355.tex
% Self-contained: no external bibliography, images, or \input files.
% Numeric section labels and visible numbers are original UnsolvedMath IDs.
\documentclass[11pt,a4paper]{article}
\usepackage[margin=24mm,headheight=14pt]{geometry}
\usepackage{fontspec}
\setmainfont{Latin Modern Roman}
\setsansfont{Latin Modern Sans}
\setmonofont{Latin Modern Mono}
\usepackage{amsmath,amssymb,mathtools}
\usepackage{unicode-math}
\setmathfont{Latin Modern Math}
\usepackage{microtype}
\usepackage{enumitem,xurl,needspace}
\usepackage[dvipsnames]{xcolor}
\usepackage{hyperref}
\hypersetup{unicode=true,colorlinks=true,linkcolor=MidnightBlue,urlcolor=MidnightBlue,
 pdftitle={UnsolvedMath Problem 1355},pdfauthor={Source-annotated compilation},bookmarksnumbered=true}
\usepackage{bookmark,tocloft,titlesec,fancyhdr}
\setlength{\cftsecnumwidth}{4.2em}
\cftsetpnumwidth{2.5em}
\cftsetrmarg{4em}
\titleformat{\section}{\Large\bfseries\raggedright}{\thesection}{0.7em}{}
\pagestyle{fancy}\fancyhf{}
\fancyhead[L]{\small\sffamily Additional UnsolvedMath statements}
\fancyhead[R]{\small Original catalogue IDs}
\fancyfoot[C]{\thepage}
\setlength{\parindent}{0pt}
\setlength{\parskip}{5pt plus 1pt}
\setlength{\emergencystretch}{3em}
\setcounter{tocdepth}{1}\setcounter{secnumdepth}{1}
\setlist[itemize]{leftmargin=3em,itemsep=4pt,topsep=4pt,parsep=0pt}
\allowdisplaybreaks
\newcommand{\N}{\mathbb{N}}
\newcommand{\Z}{\mathbb{Z}}
\newcommand{\Q}{\mathbb{Q}}
\newcommand{\R}{\mathbb{R}}
\newcommand{\C}{\mathbb{C}}
\newcommand{\F}{\mathbb{F}}
\newcommand{\A}{\mathbb{A}}
\newcommand{\PP}{\mathbb{P}}
\newcommand{\E}{\mathbb{E}}
\newcommand{\eps}{\varepsilon}
\newcommand{\dd}{\,\mathrm d}
\newcommand{\sref}[2]{\hyperlink{source:#1:#2}{[S#2]}}
\newif\ifshowcatalogue
\showcataloguetrue
\newcommand{\cataloguescope}[1]{\ifshowcatalogue
 \paragraph{Statement recorded in the supplied source.}
 #1\fi}
\begin{document}
% UnsolvedMath ID 1355; source catalogue code ST-009

\setcounter{section}{1354}

\section{A Jónsson Algebra of Cardinality Aleph-Omega}\label{1355}

{\small\sffamily UnsolvedMath ID 1355 · Set Theory}

\subsection{Definitions and mathematical statement}

\paragraph{Shared notation.}Unless a section says otherwise, $\N=\{1,2,\ldots\}$,
$\Z$ is the ring of integers, $\Q,\R,\C$ are the rational, real, and complex
fields, $[N]=\{1,\ldots,\lfloor N\rfloor\}$, and $\log$ is the natural
logarithm. For real functions of a parameter tending to infinity,
$f\ll g$ and $f=O(g)$ mean $|f|\leq Cg$ eventually for some fixed $C$;
$f\gg g$ reverses this comparison; $f=o(g)$ means $f/g\to0$;
$f\sim g$ means $f/g\to1$; and $f\asymp g$ means both order bounds.
A subscript on $O$ or $\ll$ specifies allowed constant dependence.
The expression $n^{o(1)}$ in an upper bound means growth smaller than
$n^\epsilon$ for each fixed $\epsilon>0$. Local definitions govern any
exception, finite-field convention, or use of the same symbol for another
object. An asymptotic claim is not automatically an effective algorithm.



Use an algebra with a countable collection of finitary operations. A subalgebra is a subset closed under every operation, including nullary operations when present. The algebra is Jónsson if it is infinite and every proper subalgebra has strictly smaller cardinality. The underlying set is required to have cardinality $\aleph_\omega=\sup_{n<\omega}\aleph_n$; allowing a separate constant for every element would trivialize the question. \sref{1355}{1} \sref{1355}{2}

\paragraph{Statement recorded in the supplied source.}
Does there exist a Jónsson algebra on $\aleph_\omega$? \sref{1355}{1}

\subsection{Short English statement}

Can a countably specified algebra of size aleph-omega have no proper subalgebra of the same size?

\subsection{Research attempt: term closures at a singular cardinal and an exact generation criterion}
\textbf{Status.} No Jónsson algebra of size $\aleph_\omega$ is constructed. The primary 2025 slides [R1] explicitly distinguish the longstanding problem at $\aleph_\omega$ from nearby successor-cardinal questions. The imported title about $\aleph_{\omega+1}$ in [S2] concerns a different cardinal and cannot answer this target by renaming its universe.

\subsubsection{1. The closure-size bound}
Let $\mathcal A$ have a countable finitary signature. Starting from $X\subseteq A$, put $X_0=X$ together with the values of all constants, and define $X_{n+1}$ by adjoining all outputs of basic operations on finite tuples from $X_n$. Then
\[\langle X\rangle_{\mathcal A}=\bigcup_{n<\omega}X_n.
\]
The union is closed because each operation uses only finitely many arguments, all present together at some finite stage. It is contained in every subalgebra containing $X$ by induction, so it is exactly the generated subalgebra. Countability of the signature and finitariness give
\[|\langle X\rangle_{\mathcal A}|\leq\max(|X|,\aleph_0).\tag{1}\]
For finite $X$ the right side remains countable; for infinite $X$ this uses the ZFC identities $|X|^n=|X|$ for finite positive $n$ and a countable union bound. It does not apply to operations of infinite arity or an uncountable supply of constants.

\subsubsection{2. Every such algebra is a countable union of small subalgebras}
Identify the universe with $\kappa=\aleph_\omega$ and choose increasing $X_n$ of cardinality $\aleph_n$ whose union is $\kappa$. Let $A_n=\langle X_n\rangle$. Equation (1) gives $|A_n|=\aleph_n<\kappa$, and the subalgebras are increasing with union $A$. Thus every algebra in the question, including a hypothetical Jónsson algebra, is a union of a countable chain of proper subalgebras.

There is no contradiction: Jónsson forbids a \emph{proper subalgebra of size $\kappa$}, not a union of smaller proper subalgebras equal to the whole algebra. The cofinal sequence of smaller cardinalities is exactly why the countable union can have size $\kappa$. An argument replacing that union by a uniformly bounded cardinal would misuse singularity.

\subsubsection{3. A precise criterion in terms of terms}
The Jónsson property is equivalent to
\[X\subseteq A, |X|=\kappa\quad\Longrightarrow\quad\langle X\rangle=A.\tag{2}\]
If $A$ is Jónsson, the generated subalgebra in (2) has size $\kappa$ and therefore cannot be proper. Conversely a proper subalgebra $B$ of size $\kappa$ would contradict (2) by taking $X=B$.

There are countably many terms in a countable finitary signature. Enumerate their induced operations, allowing finite arities $r_j$, as $t_j^\mathcal A:A^{r_j}\to A$. Then (2) is exactly
\[\bigcup_{j<\omega}t_j^\mathcal A[X^{r_j}]=A
\quad\text{for every }X\in[A]^\kappa.\tag{3}\]
This reformulation includes constants and repeated variables. Omitting composite terms would give only one-step closure and would not characterize generation.

A sufficient combinatorial construction would be a function $f:[\kappa]^{<\omega}\to\kappa$ with $f``[X]^{<\omega}=\kappa$ for every $X\in[\kappa]^\kappa$. Give the universe one operation of arity $n$ for each $n<\omega$, returning $f$ of the set of its arguments. Any size-$\kappa$ closed subset would contain that entire image and equal $\kappa$. This is a proved implication, not a claim that such a function has been produced.

\subsubsection{4. Verification and remaining gap}
The proof audit checks the countable signature, finite arity, inclusion of nullary operations, and use of terms rather than just basic operations. The small-cardinal example $(\mathbb N,0,s)$, where $s(n)=n+1$, has no proper subalgebra at all, but (1) shows why a countable set of generators cannot trivially create the uncountable example requested. Neither the chain decomposition nor the equivalent generation criterion forces or precludes an algebra satisfying (3). Constructing the required operations uniformly against every size-$\kappa$ subset, or showing every proposed countable signature admits a proper size-$\kappa$ closed subset, is still missing.

\subsection{Sources}

{\small

\begin{itemize}

\item[\hypertarget{source:1355:1}{S1}] ULAM.AI, \emph{UnsolvedMath}, supplied \texttt{problems.json}, record 1355 (ST-009), statement and background. The embedded literature-review date is 2026-08-17; that is the archive’s review date, not a fresh certification. Dataset landing page: \url{https://huggingface.co/datasets/ulamai/UnsolvedMath}.

\item[\hypertarget{source:1355:2}{S2}] S. Shelah, aleph\_(omega+1) has a Jónsson algebra. \url{https://doi.org/10.1007/BF02767195}. Bibliographic information imported from the supplied record.

\item[\hypertarget{source:1355:3}{S3}] UnsolvedMath, ST-009 problem record. \url{https://www.unsolvedmath.com/problems?category=10&status=open}. Bibliographic information imported from the supplied record.

\item[R1] M. Golshani, Questions that haunt my nights, while Shelah might answer by dawn, 2025 slides, discussion of Jónsson cardinals at slide 15.
\url{https://www.dmg.tuwien.ac.at/fb8/2025_slides/Golshani.pdf}.
\end{itemize}

}

\label{end:1355}
\end{document}
